\documentclass[a4paper,12pt]{article}
\usepackage[T1]{fontenc} 

%%%%%%%%%%%%%%%%%%%%%%%%%%%%%%%% CONSTANTES %%%%%%%%%%%%%%%%%%%%%%%%%%%%%%%%%%%
\newcommand{\numero}{4}                                    %Numéro de la série -1

\usepackage[french]{babel}
\usepackage[utf8]{inputenc}
\usepackage{answers}

\usepackage{hyperref}
\usepackage{multicol}

\usepackage[table,xcdraw]{xcolor}
\usepackage{listings}
\definecolor{ForestGreen}{RGB}{34,139,34}


\usepackage{enumitem}

\AtBeginDocument{\def\labelitemi{$\bullet$}}


\newcommand{\py}{\lstinline{Python} }


\definecolor{backcolour}{rgb}{0.95,0.95,0.92}

\lstset{%
	language         = Python,
	backgroundcolor  = \color{backcolour},
	basicstyle       = \ttfamily, % \upshape\ttfamily,
	keywordstyle     = \bfseries\color{blue}, %\bfseries,
	stringstyle      = \color{magenta},
	commentstyle     = \color{ForestGreen},
	alsoletter = > ,
	morekeywords = {>>>,as,assert,False,None, nonlocal,True, with,yield , <<, >>, :},
	showstringspaces = false,
	numbers=left,
	stepnumber=1,
	literate={à}{{\`{a}}}1 {é}{{\'e}}1 {è}{{\`{e}}}1 {ê}{{\^{e}}}1 {Ê}{{\^{E}}}1 {î}{{\^i}}1 {ô}{{\^{o}}}1 {ç}{{\c{c}}}1 {Ç}{{\c{C}}}1
}

\newcommand{\itemb}[1]{\item \textbf{#1}}

\usepackage{fancyhdr}  %package pour en-tetes et pied de pages
\usepackage{sectsty} % Permet de faire des modifications de police dans diverses sections des "headings" (cf. modif presentation de la page)
\pagestyle{fancy}       %Style pour en-tetes et pieds de pages
\fancyhead[CO,CE]{\sc Série 1\hspace{0.5mm}}
\fancyhead[RO,LE]{Collège Sismondi}  % LaTeX/TEX define \strut to be an invisible box of width zero that extends just enough above and below the baseline. Cela permet d'augementer légèrement la taille en bas de la box de manière à ce qu'elle soit collée à la ligne.
\fancyhead[LO,RE]{\small\ \textsl{1\textsuperscript{ère} année - DO Informatique}}
\fancyfoot[RO,LE]{2021 - 2022}
\fancyfoot[LO,RE]{\small }
\fancyfoot[CO,CE]{\thepage}

\fancyhfoffset[l]{1.2cm} % le "l" en paramètre permet d'indiquer qu'on ne veut modifier que la marge à gauche.
\renewcommand{\headrule}{{%
		\hrule \headwidth \headrulewidth \vskip-\headrulewidth}}
\renewcommand\footrulewidth{\headrulewidth}
\renewcommand{\footrule}{{%
		\vskip-\footruleskip\vskip-\footrulewidth
		\hrule \headwidth \footrulewidth\vskip\footruleskip}}

\usepackage{tikz}
%-------------------------------------------------------------------------------
%---- Eclairage : en encadré sur fond jaune avec symbôle "ampoule" à gauche ----
%-------------------------------------------------------------------------------
\definecolor{coleclairage}{RGB}{255 , 221 , 156}
\definecolor{contoureclairage}{RGB}{255 , 192 , 0}
\newenvironment{eclairage}
{
	\begin{center}%
		\begin{tikzpicture}%
			\node[rectangle, draw=contoureclairage, top color=coleclairage!50, bottom color=coleclairage!140, rounded corners=5pt, inner xsep=5pt, inner ysep=6pt, outer ysep=10pt]\bgroup                     
			\begin{minipage}{0.98\linewidth}
				\begin{minipage}{0.08\linewidth}\centerline{\includegraphics[scale=1]{Symbole_eclairage.png}}\end{minipage}
				\begin{minipage}{0.89\linewidth}\itshape\footnotesize
				}
				{                		
				\end{minipage}
			\end{minipage}\egroup;%
		\end{tikzpicture}%
	\end{center}%
}

%-------------------------------------------------------------------------------
%---- apprendre : en encadré sur fond jaune avec symbôle "ampoule" à gauche ----
%-------------------------------------------------------------------------------
\definecolor{colapprendre}{RGB}{50,205,50}
\definecolor{contourapprendre}{RGB}{34,139,34}
\newenvironment{apprendre}
{
	\begin{center}%
		\begin{tikzpicture}%
			\node[rectangle, draw=contourapprendre, top color=colapprendre!10, bottom color=colapprendre!50, rounded corners=5pt, inner xsep=5pt, inner ysep=6pt, outer ysep=10pt]\bgroup                     
			\begin{minipage}{0.98\linewidth}
				\begin{minipage}{0.08\linewidth}\centerline{\includegraphics[width=30px]{Symbole_learn.png}}\end{minipage}
				\begin{minipage}{0.89\linewidth}\itshape\footnotesize
				}
				{                		
				\end{minipage}
			\end{minipage}\egroup;%
		\end{tikzpicture}%
	\end{center}%
}

\definecolor{colimportant}{RGB}{247 , 189 , 164}
\definecolor{contourimportant}{RGB}{237 , 125 , 49}
\newenvironment{important}
{
	\begin{center}%
		\begin{tikzpicture}%
			\node[rectangle, draw=contourimportant, top color=colimportant!50, bottom color=colimportant!140, rounded corners=5pt, inner xsep=5pt, inner ysep=6pt, outer ysep=10pt]\bgroup                     
			\begin{minipage}{0.08\linewidth}\centerline{\includegraphics[scale=0.8]{Symbole_attention.png}}\end{minipage}
			\begin{minipage}{0.89\linewidth}
			}
			{                		
			\end{minipage}\egroup;
		\end{tikzpicture}%
	\end{center}%
}

%-----------------------------------------------------------------
%---- Modification présentation de la page: marges de la page ----
%-----------------------------------------------------------------
%\addtolength{\hoffset}{-1in}              % 1
%\addtolength{\voffset}{-1in}              % 2
\addtolength{\oddsidemargin}{-0.1 in} % 3
\addtolength{\evensidemargin}{-1in} % 3
\addtolength{\topmargin}{-1in}       % 4
\addtolength{\headheight}{6pt}       % 5
%\addtolength{\headsep}{-0.2cm}           % 6
\setlength{\textheight}{26cm}    % 7
\setlength{\textwidth}{16.5cm}      % 8
\addtolength{\marginparsep}{0pt}      % 9
\setlength{\marginparwidth}{0pt}   % 10
\addtolength{\footskip}{-1mm}           %11

\setlength{\parindent}{0em}% pas d'indentation


% Customiser le nom des sections
\usepackage{titlesec}
\titleformat{\section}[hang]{\Large \bfseries}{Série \thesection:\ }{0pt}{}

\renewcommand{\familydefault}{\sfdefault} % pour avoir des polices san serif

\newtheorem{Exc}{Exercice}
\Newassociation{correction}{Soln}{mycor}
\renewcommand{\Solnlabel}[1]{\bfseries Ex #1 }
\def\exo#1{%
	\futurelet\testchar\MaybeOptArgmyexoo}
\def\MaybeOptArgmyexoo{
	\ifx[\testchar \let\next\OptArgmyexoo
	\else \let\next\NoOptArgmyexoo \fi \next}
\def\OptArgmyexoo[#1]{%
	\begin{Exc}[#1]\normalfont}
	\def\NoOptArgmyexoo{%
		\begin{Exc}\normalfont}
		\newcommand{\finexo}{\end{Exc} \vspace{3mm}}
	\newcommand{\flag}[1]{}
	\newcommand{\entete}[1]

\newcommand{\getexocompteur}{{\the\numexpr \arabic{Exc}  \relax}}	
	
\newcommand{\eexo}{\vspace{5mm}} % espace pour séparer les exercices
\definecolor{light-gray}{gray}{0.95} 
\lstset{language=Python, upquote=true, showstringspaces=false, columns=fullflexible, basicstyle = \ttfamily, backgroundcolor = \color{light-gray}, numbers=right, stepnumber=1, keywordstyle=\color{blue}, stringstyle=\color{magenta}, commentstyle=\color{red}, breaklines=true,} 

\begin{document}
%		\title{\vspace{-3cm}Série 1}
%		\date{\vspace{-2cm}}
%		\maketitle

\fancyhead[CO]{\sc Série \arabic{section} \hspace{0.5mm}}
%\fancyhead[CE]{\sc Série \arabic{section} \hspace{0.5mm}}
\setcounter{section}{\numero}

\section{Qui cherche trouve}
\Opensolutionfile{mycor}[cor01]
%%%%%%%%%%%%%%%%%%%%%%%%%%%%%%%% EXERCICE %%%%%%%%%%%%%%%%%%%%%%%%%%%%
\exo{}[Deviner un nombre]  ~\\ 
Vous devez réaliser un programme dans lequel un joueur ou une joueuse doit trouver un nombre compris entre 1 et 100. Les étapes à réaliser sont les suivantes:

\begin{enumerate}
\item L'ordinateur choisit au hasard un nombre entre 1 et 100.
\item Le programme demande au joueur ou à la joueuse d'entrer un nombre.
\item À chaque essai l'ordinateur indique si le nombre à trouver est plus grand ou plus petit que le nombre tiré au hasard par l'ordinateur.
\end{enumerate}

Prenons un exemple où l'ordinateur a tiré 72 comme nombre aléatoire:

\begin{lstlisting}
Entrez un nombre: 50
Le nombre au hasard est plus grand que 50
Entrez un nombre: 75
Le nombre au hasard est plus petit que 75
Entrez un nombre: 70
Le nombre au hasard est plus grand que 70
Entrez un nombre: 72
Bravo vous avez trouvé !
\end{lstlisting}
\hfill
	
	\vspace{4pt}
		

Pour utiliser la fonction randint qui retourne un nombre entier au hasard, vous devez importer son module \footnote{On appelle “module” tout fichier constitué de code Python (c'est-à-dire tout fichier avec l'extension .py ) importé dans un autre fichier ou script. Les modules permettent la séparation et donc une meilleure organisation du code.} de la manière suivante dans les premières lignes de votre code:

\begin{lstlisting}
from random import *
\end{lstlisting}
Une fois importé, vous pouvez utiliser la fonction de cette manière:
\begin{lstlisting}
nombre_entier_hasard = randint(1,100)
\end{lstlisting}

\hfill
	
	\vspace{4pt}
		

Dans votre programme, créer une fonction qui aura comme paramètres le nombre au hasard et le nombre saisi par le joueur ou la joueuse. Cette fonction retourne \lstinline{True} si le nombre est trouvé sinon \lstinline{False}. Elle devra également afficher si le nombre saisi est trop grand ou trop petit. 
\begin{lstlisting}
def test_nb(nb_hasard, nb_saisi):
...
\end{lstlisting}
\hfill
	
	\vspace{4pt}
		
Pour demander un nombre au·à la joueur·euse, vous pouvez utiliser la fonction \lstinline{input}:
\begin{lstlisting}
reponse = input("Question : ") 
\end{lstlisting}
qui retourne la chaîne de caractère entrée par l'utilisateur·trice.
\begin{lstlisting}
from random import *

def test_nb(nb_hasard, nb_saisi):
      ...
      
nombre_entier_hasard = randint(1,100)
...
nombre_entier_choisi = int(input("Choisissez un nombre entre 1 et 100: "))
\end{lstlisting}

\begin{correction}
		~\\ \vspace{-5pt}
	\begin{lstlisting}
from random import *

def test_nb(nb_hasard, nb_saisi):
### Fonction qui compare 2 nombres et retourne True s'ils sont égaux, False sinon ###
    if nb_hasard > nb_saisi:
        print("Plus grand...")
    elif nb_hasard < nb_saisi:
        print("Plus petit...")
    return nb_hasard == nb_saisi

# Début de partie
nombre_entier_hasard = randint(1,100)
gagne = False
# print("Nombre au hasard:",nombre_entier_hasard)

# Boucle de jeu
while not gagne:
    nombre_entier_choisi = int(input("Choisissez un nombre entre 1 et 100: "))
    gagne = test_nb(nombre_entier_hasard, nombre_entier_choisi)

#Fin de partie
print("C'est gagné !")
\end{lstlisting}
\hfill
	\end{correction}
\finexo
%%%%%%%%%%%%%%%%%%%%%%%%%%%%%%%% EXERCICE %%%%%%%%%%%%%%%%%%%%%%%%%%%%
\exo{}[Suite]  ~\\ 
Complétez le programme de l'exercice précédent de la manière suivante (pour chaque étape, faites un nouveau programme en ajoutant le nouveau comportement) :
\begin{enumerate}[topsep=0pt,itemsep=0pt,parsep=30pt,partopsep=0pt]
\item Créer une fonction qui vérifie les deux points suivants:
    \begin{enumerate}
    \item Est-ce bien un nombre entier qui a été saisi ?
          Utiliser la méthode\footnote{Une méthode est une fonction qui "appartient à" un objet en programmation objet.} \lstinline{isnumeric()}. Exemple d'utilisation:
        \begin{lstlisting}
         # retournera True si la variable mon_nombre contient un nombre entier.
         if mon_nombre.isnumeric():
             ...
        \end{lstlisting}
    \item Est-il bien compris entre 1 et 100 ?
    \end{enumerate}
    Cette fonction devra retourner \lstinline{True} si le nombre est bien entier et compris entre 1 et 100. Sinon \lstinline{False}. Ensuite redemandez à l'utilisateur·trice de ressaisir un nombre si besoin.
\item Ajoutez un compteur qui compte le nombre de tentatives. À la fin de la partie, le programme affiche ce nombre:
\begin{lstlisting}
Vous avez gagné en 7 coup(s)
\end{lstlisting}
\item Limité le nombre de tentatives à 10. Si le·la joueur·euse n'a pas trouvé en 10 tentatives, le programme donne ensuite la réponse. Par exemple:
\begin{lstlisting}
Vous avez joué 10 fois sans trouver... La réponse était 42.
\end{lstlisting}
\item Après 5 essais, le programme donne le nombre de tentatives restantes:
\begin{lstlisting}
Vous avez encore 5 tentatives.
\end{lstlisting}
\item À la fin de la partie, le programme propose de rejouer. \\
Pour simplifier, le·la joueur·euse entre 0 ou 1:
\begin{lstlisting}
Voulez-vous rejouer? Taper 0 pour NON ou taper 1 pour OUI: 
\end{lstlisting}
\end{enumerate}

\begin{correction}
		~\\ \vspace{-5pt}
	...
\end{correction}
\finexo

% Solution		
%		\newpage
%		\setcounter{page}{1}
%		\setcounter{section}{\numero}
%		\Closesolutionfile{mycor}
%		\titleformat{\section}[hang]{\Large \bfseries}{Corrigé Série \thesection:\ }{0pt}{}
%		
%
%		\fancyhead[CO]{\sc Corrigé Série \arabic{section} \hspace{0.5mm}}
%		\section{}
%		\Readsolutionfile{mycor}
	\end{document}